\section{TitaNet}
For extracting speaker representations (embeddings), Koluguri et al. \cite{koluguri2021} proposed a novel network architecture called TitaNet.
TitaNet simply takes a variable length speech signal as input and outputs a fixed length speaker embedding.
Their architecture achieves state-of-the-art performance on both speaker verfications and speaker diarization tasks.


\subsection{Model Architecture}
The model is based on encoder and decoder structure. The encoder is used for feature exctraction and then those features are passed to attention pooling layer 
in order to calculate attention features across channels dimension.

\begin{figure}[ht]
    \centering
    \includegraphics[width=0.8\textwidth]{images/related_work/titanet/arch.png}
    \caption{TitaNet Encoder and Decoder Architecture}
    \label{fig:TitaNet}
\end{figure}

The encoder decoder architecture of TitaNet is shown in Figure \ref{fig:TitaNet}. 

\subsubsection{Encoder}
The model encoder is based on ContextNet-$B\times R \times C$ model such that:
\begin{itemize}
    \item $B$: number of blocks
    \item $R$: number of sub-blocks per block
    \item $C$: number of filters in the convolution layers of each block
\end{itemize}
The encoder starts with prologue block $B_{0}$ followed by $N-2$ mega blocks and finally it ends with epilogue block $B_{N}$.
Both prologue and epilogue blocks are different from other mega blocks; they have convolution layer followed by batch normalization and ReLU layers but they differ in kernel size 
(3 and 1 for prologue and epilogue respectively).
As for the mega blocks, they consist of time-channel separable convolution (with stride 1 and dilation 1), followed by batch normalization, ReLU and dropout.
Those layers are repeated $R$ times to allow varying the depth. For global context, these channels are residually connected with Squeeze and Excitation layer.

\subsubsection{Decoder and Embeddings}
The high level features computed by the encoder are then passed to decoder to get the speaker embedding.
Attentive statistics pooling layer is used to get a time-independent feature representation $S$ of size $B \times 3072$.
These features are then passed to 2 linear layers:
\begin{itemize}
    \item $l_1$: Output a fixed-size embedding of size 192
    \item $l_2$: Output a vector of probabilities that the given utterance belongs to one of the training set speakers
\end{itemize}

\subsubsection{Loss funciton}
TitaNet uses additive angular margin (AAM) loss. It helps to optimize the cosine distance between speaker embeddings:
\begin{equation}
    \text{L} = - \frac{1}{N} \sum_{i=1}^{N} \log \frac{e^{s(\cos(\theta_{y_i} + m))}}{e^{s(\cos(\theta_{y_i} + m))} + \sum_{j=1, j \neq y_i}^{n} e^{s \cos \theta_{j}}}
\end{equation}
where $m$ is margin, $s$ is scale and $\theta_j$ is the angle between the final layer weight $W_j$ and incoming feature $x_i$.
Both $m$ and $s$ are predefined hyper parameters.

\subsection{Experimental Setup}
Three TitaNet models were designed: TitaNet-S, TitaNet-M and TitaNet-L with 256, 512 and 1024 channels respectively.
All the mentioned models have the same $R$ as 3 and the same filter sizes as 3, 7, 11 and 15.
\subsubsection{Dataset}
The authors used VoxCeleb1, VoxCeleb2 dev, LibriSpeech as well as other datasets with a total number of speakers of 16681 and a duration of 3373 hours.
The training data is also augmented using RIR impulse corpora as well as different speed pertubation (0.95x and 1.05x).


\subsection{Conclusion}
This paper presents TitaNet, a novel model for speaker representation learning that uses global context of SE layers along with channel attentive statistics pooling for the extraction of fixed length speaker embedding.
All models' implementations and checkpoints are available in NeMo toolkit (provided by NVIDIA).    